% Source PDF pages 53--54; manual pages 2-24--2-25.
\subsection{Velocity Cartesian Coordinate Systems}\label{sec:velocity-coordinate}

Intermediate Cartesian systems, called velocity coordinate systems, are used to
convert from geocentric or geodetic coordinates to the wind coordinate system
(\cref{sec:wind-coordinate}). One velocity system is referenced to the local
geocentric horizon system and another to the local geodetic horizon system. They
are defined identically except for the reference-frame subscripts. The development
below uses the geodetic system; replacing the $gd$ subscripts with $gc$ gives the
geocentric form.

The origin of a velocity coordinate system is at the vehicle's center of mass. For
the air-relative velocity system, the $x$ axis is aligned with the wind-corrected
velocity vector
\begin{equation}
  \vec V_{w\oplus}=\vec V_{\oplus}-\Delta\vec V_{w\oplus},
  \label{eq:wind-corrected-velocity}
\end{equation}
where $\Delta\vec V_{w\oplus}$ is the input wind velocity vector relative to the
earth's surface, as described in \cref{sec:block-wind}. The vector
$\vec V_{w\oplus}$ is therefore the vehicle velocity relative to the surrounding air.
Its magnitude is the output variable \statevar{vair}, or airspeed.

The velocity-system $y$ axis is perpendicular to the plane formed by its $x$ axis
and the geodetic $\hat z_{gd}$ axis. The $z$ axis is perpendicular to the other two
and points generally toward the earth, as shown in
\cref{fig:velocity-coordinate-system}.

\begin{figure}[htbp]
  \centering
  \resizebox{0.72\linewidth}{!}{\begin{tikzpicture}[scale=0.88,>=Latex,line cap=round,line join=round]
  \coordinate (O) at (0,0);
  \coordinate (P) at (4.15,2.60);
  \draw[->] (O)--(0,4.25) node[above] {$\hat{z}_{\oplus}$};
  \draw[->] (O)--(5.35,0) node[right] {Equatorial plane};
  \draw[line width=0.9pt] (0,3.35) arc[start angle=90,end angle=5,x radius=4.55,y radius=3.35];
  \draw[->,line width=1pt] (O)--(P) node[midway,above] {$\vec r$};
  \fill (P) circle (1.6pt) node[right=7pt] {Vehicle};
  \draw[->,line width=1.1pt] (P)--($(P)+(-1.65,0.88)$) node[above left] {$\hat{x}_{V_w}$};
  \draw[->,line width=1.1pt] (P)--($(P)+(0.10,-1.20)$) node[below] {$\hat{z}_{V_w}$};
  \draw[->,line width=1.1pt] (P)--($(P)+(-0.95,-0.28)$) node[below left] {$\hat{y}_{V_w}$};
  \draw[->,line width=1.1pt] (P)--($(P)+(-1.30,1.15)$) node[above left] {$\vec V_{w\oplus}$};
  \draw[->,line width=0.9pt] (P)--($(P)+(-0.55,-0.80)$) node[below left] {$\hat{z}_{gd}$};
  \node[align=left] at (5.55,3.35) {Note: $\hat{y}_{V_w}$ is\\into the page};
\end{tikzpicture}
}
  \caption{Velocity coordinate system.}
  \label{fig:velocity-coordinate-system}
\end{figure}

The function \texttt{wind\_unit\_vectors} computes the air-relative velocity-system
unit vectors from
\begin{equation}
  \hat x_{V_w}=\frac{\vec V_{w\oplus}}
  {\lVert\vec V_{w\oplus}\rVert},
  \label{eq:air-relative-velocity-x-axis}
\end{equation}
\begin{equation}
  \hat y_{V_w}
  =\frac{\hat z_{gd}\times\hat x_{V_w}}
  {\lVert\hat z_{gd}\times\hat x_{V_w}\rVert},
  \label{eq:air-relative-velocity-y-axis}
\end{equation}
\begin{equation}
  \hat z_{V_w}
  =\frac{\hat x_{V_w}\times\hat y_{V_w}}
  {\lVert\hat x_{V_w}\times\hat y_{V_w}\rVert}.
  \label{eq:air-relative-velocity-z-axis}
\end{equation}

A second, analogous velocity system relates the earth-relative velocity vector to
the geodetic system through the flight-path angles. Its $x$ axis is aligned with
$\vec V_{\oplus}$ rather than with $\vec V_{w\oplus}$:
\begin{equation}
  \hat x_V=\frac{\vec V_{\oplus}}
  {\lVert\vec V_{\oplus}\rVert}.
  \label{eq:earth-relative-velocity-x-axis}
\end{equation}
The geodetic system is related to this velocity system by a rotation about
$\hat z_{gd}$ through the heading angle $\psi_{gd}$, followed by a rotation about
the intermediate $y$ axis through the vertical flight-path angle $\gamma_{gd}$:
\begin{equation}
  \begin{bmatrix}
    \hat x_V\\[0.2em]
    \hat y_V\\[0.2em]
    \hat z_V
  \end{bmatrix}
  =
  \TaosPitchMatrix{\gamma_{gd}}
  \TaosYawMatrix{\psi_{gd}}
  \begin{bmatrix}
    \hat x_{gd}\\[0.2em]
    \hat y_{gd}\\[0.2em]
    \hat z_{gd}
  \end{bmatrix}.
  \label{eq:geodetic-to-velocity-transformation}
\end{equation}
Expanding the matrices gives
\begin{equation}
  \hat x_V
  =\cos\gamma_{gd}\cos\psi_{gd}\,\hat x_{gd}
  +\cos\gamma_{gd}\sin\psi_{gd}\,\hat y_{gd}
  -\sin\gamma_{gd}\,\hat z_{gd},
  \label{eq:velocity-x-from-geodetic}
\end{equation}
\begin{equation}
  \hat y_V
  =-\sin\psi_{gd}\,\hat x_{gd}
  +\cos\psi_{gd}\,\hat y_{gd},
  \label{eq:velocity-y-from-geodetic}
\end{equation}
\begin{equation}
  \hat z_V
  =\sin\gamma_{gd}\cos\psi_{gd}\,\hat x_{gd}
  +\sin\gamma_{gd}\sin\psi_{gd}\,\hat y_{gd}
  +\cos\gamma_{gd}\,\hat z_{gd}.
  \label{eq:velocity-z-from-geodetic}
\end{equation}
These relationships are used in the derivation of the flight-path-angle rates in
\cref{sec:flight-path-rates}.
